\section{Part VI: Ownership, Dilution, Terms and Exit Waterfalls}

\subsection{116. 'How much of the company do I own?' requires a denominator}


The ownership fraction depends on shares outstanding, outstanding convertible instruments, options promised or issued, warrants, unallocated option pool and the financing definition of \textbf{fully diluted capitalization}. A cap table should distinguish legal issued shares from an economic pro forma view in which relevant conversion and vesting events occur. A founder who owns 90\% of issued common shares may have substantially less on a fully diluted and post-financing basis.

Keep a real security-by-security ledger: holder, class, certificates, grant date, vesting, strike price, conversion rights, cap and discount, protective provisions, proceeds rights and transfer restrictions. Shares must reconcile with legally authorized corporate actions. A simplified spreadsheet is useful; undocumented promises of equity are not.

\subsection{117. Pre-money and post-money valuation arithmetic}


In the simplest priced-equity financing with no complicating instruments or option pool changes: $\text{Post-Money} = \text{Pre-Money} + \text{New Investment}$. If a company has a ₹12 crore pre-money valuation and raises ₹3 crore, post-money is ₹15 crore and the new investor buys \textbf{20\%}. Existing shareholders collectively fall from 100\% to 80\% of the \textit{larger} equity base. If the founder held 80\% before the round, afterward the founder holds \textbf{64\%} absent further changes.

This is capitalization arithmetic, not a proof that the company's operating assets are now worth ₹15 crore in cash. Negotiated valuations may include strategic bargaining and securities preferences. Secondary sales do not inject equivalent primary operating cash into the company. Distinguish whose pocket receives the money.

\subsection{118. A multi-round dilution example}


Assume the founder initially owns 100\% of common shares. Round A sells \textbf{20\% post-money} to new investors, leaving founder \textbf{80\%}. Later Round B sells \textbf{25\% post-money} to additional investors, with no other changes. The founder becomes $80\% \times 75\% = 60\%$, Round A becomes \textbf{20\%$\times$75\% = 15\%}, and Round B holds \textbf{25\%}. The cap table totals 100\%.

If a 10\% post-round option pool is then freshly created by issuing shares to employees, everyone else is diluted proportionally: founder \textbf{54\%}, Round A \textbf{13.5\%}, Round B \textbf{22.5\%}, pool \textbf{10\%}. Real financing documents often require a pool expansion \textit{before} the priced round, shifting more dilution onto existing owners. The \textbf{order of events} is a negotiated economic term.

\subsection{119. Option pools: compensation is real dilution}


Employee equity supports recruitment and alignment, especially where cash salaries are constrained or the company might appreciate substantially. But options may have exercise price, vesting, tax, expiry, liquidity, repurchase, and termination conditions. An 'employee pool of 12\%' is not equivalent to 12\% already earned by employees. Founders must distinguish \textbf{issued grants}, \textbf{vested options}, \textbf{exercised shares}, and \textbf{unallocated capacity}.

The economic cost of employee equity is partly dilution and partly the opportunity to attract talent who improve enterprise value. A grant can be excellent if it helps recruit a high-impact engineer or manager and avoids unsustainable cash compensation. It is wasteful when distributed indiscriminately or explained as guaranteed future wealth. Give candidates understandable documents and avoid presenting gross notional option value as liquid salary.

\subsection{120. The option-pool shuffle}


Suppose an investor offers a certain headline pre-money valuation but requires a new option pool to be included in the pre-money capitalization. The pool effectively reduces the pre-round owners' percentage \textbf{before} the investor's money arrives. As a result, the founder may retain less than the simple investment/post-money calculation suggests.

Never compare term sheets only on 'pre-money valuation'. Build post-closing cap tables for \textbf{all outstanding and promised claims}, including the new option pool, converting securities, warrants and fees paid in equity. Ask who bears each portion of dilution. 2 economically different offers can produce the same press-release valuation.

\subsection{121. Post-money SAFEs and multiple caps}


YC's simplified post-money SAFE logic makes aggregate potential dilution more legible: a hypothetical \textbf{₹20 lakh SAFE at a ₹4 crore post-money cap} corresponds to approximately \textbf{5\%} of the defined pre-new-money capitalization under the assumed conversion terms. Another \textbf{₹20 lakh at the same cap} adds approximately another 5\%, creating a 10\% aggregate claim before a new priced round. Actual treatment depends on instrument terms, capitalization definitions and conversion conditions. The example is arithmetic, \textbf{not advice to issue such an instrument in India}.

A founder who signs 6 'small' post-money SAFEs can sell a surprisingly large portion of the company before receiving a priced-round offer. Track every investment amount, valuation cap, discount, side letter and pro rata right in a fully diluted model updated after every signature.

\subsection{122. Cap and discount conversion scenarios}


A convertible instrument may permit conversion at a valuation cap, a discount to the next round or whichever yields the investor a more favorable price. The resulting share count depends on the defined capitalization and the next financing price. A cap has its strongest effect when the company raises at a price well above the cap; a discount may govern when the next price is lower. Interest accrual on notes can increase converting principal.

Model at least 3 next-round prices, an adverse down round and \textbf{no qualified financing before maturity} for notes. A cap is not a guaranteed company valuation. Preferred conversion may come with rights that differ from those of ordinary shareholders. Do not mix the ownership approximation for a post-money SAFE with the mechanics of a pre-money note as if they were identical.

\subsection{123. Liquidation preference: seniority matters}


A typical \textbf{1$\times$ non-participating liquidation preference} gives preferred investors a choice-subject to specific terms-between receiving their invested amount before common holders or converting into common and sharing pro rata, whichever is economically better. Suppose an investor invests \textbf{₹2 crore for 20\%}. At a \textbf{₹5 crore} distributable exit with no debt or transaction costs, the investor would normally prefer \textbf{₹2 crore preference} to 20\% of ₹5 crore = ₹1 crore. Common holders receive ₹3 crore. At a \textbf{₹20 crore} exit, conversion to 20\% yields \textbf{₹4 crore}, larger than the ₹2 crore preference.

This simple example assumes a single series, 1$\times$ non-participating preference and no other securities. Multi-round priority, dividends, conversion, seniority and escrow alter the waterfall. The headline 20\% ownership does not ensure the investor receives only 20\% of every possible exit.

\subsection{124. Participating preferred: taking the preference and the upside}


A \textbf{1$\times$ participating preferred} security may 1st recover the investment and then also share the remaining distributable proceeds at its participation percentage, often subject to a cap. Using the same \textbf{₹2 crore invested for 20\%} and a \textbf{₹5 crore exit}, the investor receives $\textnormal{₹}2\text{ cr preference} + 20\% \times \textnormal{₹}3\text{ cr} = \textnormal{₹}2.6\text{ crore}$ if uncapped and as modeled; common holders receive ₹2.4 crore. A non-participating preferred investor received ₹2 crore in the earlier case.

Participation makes a major difference in moderate exits, where founders often hope to preserve meaningful proceeds. Some term sheets offer a higher headline valuation in exchange for participation or a multiple preference. Compare distributions across low, mid and high exits before declaring the offer better.

\subsection{125. Multiple liquidation preferences and stacked rounds}


A \textbf{2$\times$ preference} can entitle investors to receive twice invested capital before common shares receive residual value, subject to terms and available proceeds. Different series can rank senior, pari passu or junior, and a later 'senior' round can change the economics of earlier investors as well as common owners. Accrued dividends can increase claims. A down round can add preferences while reducing nominal share price.

Model the exit distribution at values where the company might realistically sell, not only a visionary 1,000,000,000-dollar exit. A founder with 40\% nominal ownership can receive little in an outcome where senior claims absorb the proceeds. Preference is not inherently abusive; it is a risk allocation. \textbf{Price that allocation explicitly.}

\subsection{126. Anti-dilution protection is not the same as ordinary dilution}


Ordinary dilution occurs when new shares are issued and existing ownership percentages shrink. Contractual \textbf{anti-dilution adjustment} can change an investor's conversion price when a future round occurs below an earlier price. \textbf{Full-ratchet} protection can be particularly severe because it may reset the protected investor's conversion price to the new low share price regardless of issue size. \textbf{Broad-based weighted-average} formulas can be less punitive but still shift dilution to other holders.

Definitions, exemptions (such as certain employee grants), capitalization bases and pay-to-play provisions vary. Do not rely on a generic cap-table example for a signed agreement. Model a small down round and a large down round separately. A founder should understand whether a 1-rupee financing at a low price could trigger disproportionate changes in ownership.

\subsection{127. Pro rata, preemption and future financing access}


A pro rata or preemptive right allows a holder, within terms, to invest in later financings to preserve its ownership percentage. This can support continuity with aligned investors but may complicate allocation to new strategic investors. Some contracts have \textbf{super pro rata}, rights of 1st offer or other privileges. Founders need to distinguish the right to \textbf{invest more cash} from a free anti-dilution right: pro rata participation usually requires additional money.

A highly fragmented shareholder base with strong participation rights can slow future rounds. Consider how pro rata rights work under tight allocation, insider-led bridges and distress. Track exercise periods, notice obligations and waivers. An early side letter that seems harmless may matter greatly when the company needs rapid financing later.

\subsection{128. Control rights versus economic ownership}


Owning 60\% of economic shares does not necessarily grant unilateral control if a shareholder agreement gives minority investors veto rights over budgets, new securities, borrowing, related-party transactions, key appointments, sale of assets or changes in business. Board composition, quorum, founder vesting and protective provisions can also matter. Conversely, a founder with less than 50\% can retain influence through governance structures, though law, fiduciary duties and agreements constrain that influence.

Distinguish \textbf{ordinary-course decisions} from \textbf{reserved matters}. Investors reasonably seek safeguards against misappropriation, uncontrolled dilution and extreme strategic shifts. But a veto over routine teacher hiring or normal product changes can paralyze a small company. Negotiate clear thresholds, time limits, materiality definitions and deadlock resolution appropriate to stage and mission.

\subsection{129. Founder vesting and reverse vesting}


Vesting arrangements may cause founders or employees to earn continuing rights over time, protecting the company if someone leaves early. For founders, \textbf{reverse vesting} often means issued or granted equity remains subject to repurchase or forfeiture conditions under applicable law. Good-leaver/bad-leaver definitions, acceleration on acquisition, termination rights and pricing can determine whether a departure becomes disastrous.

The balance is legitimate: new capital may fund future founder labor, not merely historical invention. But a founder should not sign a structure that allows arbitrary termination and near-total forfeiture without understanding legal protections and governance. Founder vesting is a contractual risk allocation; it deserves the same stress testing as dilution and liquidation preference.

\subsection{130. Secondary liquidity and founder cash-outs}


A \textbf{primary financing} issues new securities and brings operating capital into the business. A \textbf{secondary transaction} transfers existing securities and pays the selling holder, not the company. A mixed round may do both. Founder secondary sales can diversify personal risk and reduce pressure for short-term salary, but investors may worry about alignment, and transaction restrictions or taxation may apply.

If a fund announces a ₹50 crore investment and ₹20 crore is founder secondary, the business did not receive ₹50 crore in its bank account. A financing plan that ignores the primary-versus-secondary split overstates runway. Analyze how any secondary transaction affects incentives, future ownership, cap-table complexity and governance.

\subsection{131. Exit rights, drag-along, tag-along and transfer restrictions}


\textbf{Drag-along rights} can require shareholders to participate in an approved sale under specified protections. \textbf{Tag-along rights} may allow minority holders to sell alongside a majority seller. Rights of 1st refusal, co-sale, transfer restrictions, lockups and redemption rights affect liquidity and bargaining. Investor put or redemption provisions can be particularly consequential if enforceable and timed before the company has available cash.

These rights are not identical across jurisdictions. An Indian private company's permissible share transfers and foreign-investment pricing conditions may constrain how some contractual rights operate. Founders should understand \textbf{what happens when the company does not sell on the expected timetable}-a much more common state than a neat venture-case exit.

\subsection{132. The venture return waterfall is not the founder's net proceeds}


A hypothetical company sells for ₹50 crore. Before calculating common shareholder proceeds, subtract transaction fees, financial debt, unpaid obligations, and relevant purchase-price adjustments; apply security seniority, preference and conversion rights; then consider escrow, earn-out and tax. The proceeds attributable to a founder's ordinary shares can differ materially from the multiplication of exit headline by nominal percentage.

A buyer may demand that founders roll over equity or continue employment to earn deferred consideration. A sale can also trigger employee option exercise or payout, indemnification and working-capital true-ups. Prepare \textbf{cash-at-close, deferred, contingent and tax-adjusted} outcomes instead of presenting a single gross exit number as personal wealth.

\subsection{133. The founder's cap-table stress test}


Maintain a capitalization model for \textbf{no next round}, \textbf{flat round}, \textbf{up round}, \textbf{down round}, \textbf{large option-pool expansion}, \textbf{bridge note}, \textbf{secondary liquidity}, and exits below/near/above total invested capital. Record each holder's rights, not only percentages. Compare capital offers under the same assumptions and verify signed documents against the spreadsheet.

Do not use an investor's preferred valuation example as the only model. Build a case where the company becomes a stable, profitable business without ever becoming a unicorn. Ask whether the financing instrument remains tolerable in that state. A round whose economics only make sense in a spectacular outcome may be poorly matched to a genuine but bounded institutional opportunity.

\subsection{134. Is dilution bad?}


Dilution is rational if the capital increases the value of the remaining stake enough, after risk and control costs. Owning \textbf{60\% of a company worth ₹100 crore} may be preferable to owning 100\% of 1 worth ₹15 crore, provided the valuations correspond to realizable economic claims and the founder is comfortable with the governance and risk. Conversely, raising money at a high paper valuation can leave a founder with smaller control, repeated financing dependence and limited proceeds under preference stacks.

The right objective is not maximum founder percentage. It is \textbf{durable founder value, mission fidelity, decision freedom and downside survivability}. A startup can use equity as a productive resource without treating dilution as failure; a bootstrap founder can preserve ownership without treating all external money as corrupting.

\subsection{135. The cost of founder dilution is nonlinear}


A reduction from 100\% to 90\% ownership can be economically different from 51\% to 41\% if it changes control, veto rights, governance coalitions or future exit flexibility. Some founder objectives require majority approval; others require special resolutions or consent from particular share classes. There may be restrictions on future related-party transactions or ability to move assets into another entity. Thus the \textit{marginal} cost of each percentage point is not constant.

Before financing, map decision rights after \textbf{this round and the next 2 plausible rounds}, not only at closing. If a venture requires 4 future rounds to fulfill the plan, preserving all control through round 1 is insufficient. Conversely, a private company with no anticipated next financing should not accept permanent restrictions designed for a rapidly refinancing venture.

\subsection{136. The best term sheet is the 1 the company can live with}


Evaluate financing on \textbf{capital received, expected dilution, economic preferences, future fundraising flexibility, downside remedies, board control, investor conduct, exit alignment and administrative burden}. Term sheets are multidimensional. Sometimes a smaller round and slower growth are better; sometimes more dilution buys a stronger company; sometimes a strategic partnership is better than investment.

The founder must understand \textbf{what the money does in a good case, an average case and a bad case}. A financing agreement is a strategic constitution for the business, not just an amount and a valuation.

